\documentclass[11pt]{article}
\usepackage[margin=1in]{geometry}
\usepackage{amsmath,amssymb,amsthm,mathtools}
\usepackage[T1]{fontenc}
\usepackage{lmodern}
\usepackage{microtype}
\setlength{\emergencystretch}{2em}
\usepackage[colorlinks=true,linkcolor=blue,citecolor=blue,urlcolor=blue]{hyperref}
\newtheorem{theorem}{Theorem}
\newtheorem{lemma}[theorem]{Lemma}
\newtheorem{proposition}[theorem]{Proposition}
\newtheorem{corollary}[theorem]{Corollary}
\theoremstyle{definition}
\newtheorem{definition}[theorem]{Definition}
\newtheorem{remark}[theorem]{Remark}
\newcommand{\ket}[1]{\lvert #1\rangle}
\newcommand{\bra}[1]{\langle #1\rvert}
\newcommand{\norm}[1]{\lVert #1\rVert}
\newcommand{\poly}{\operatorname{poly}}
\newcommand{\polylog}{\operatorname{polylog}}
\newcommand{\diag}{\operatorname{diag}}
\newcommand{\sgn}{\operatorname{sgn}}
\newcommand{\Id}{\mathrm I}
\newcommand{\F}{\mathcal F}
\newcommand{\AC}{\operatorname{AC}}
\newcommand{\RC}{\operatorname{RC}}
\newcommand{\bits}{\operatorname{bits}}
\DeclareMathOperator{\spanop}{span}
\title{A coherent near-linear interlacing transform\\and a subcubic symmetric-group QFT}
\author{Research draft from the Fourier decomposition project}
\date{22 October 2026}
\begin{document}
\maketitle

\begin{abstract}
We give a constructive proof of an approximate, clean implementation of an
integer-interlacing orthogonal matrix using
$K\poly(\log K,\log L,\log(1/\delta))$ elementary gates, including coherent
coefficient and angle computation. The method combines a logarithmic
absolute-matrix norm bound, a finite positive balancing iteration, fast
logarithmic and inverse-distance potentials, and overlap amplification.
A compact streaming implementation of the known Mackey induction recurrence
then gives an $S_n$ quantum Fourier transform with
$n^{5/2}\poly(\log(n+1),\log(1/\epsilon))$ gates. No quantum random-access
memory, precomputed diagram oracle, postselection, or exact transcendental
gate is assumed. This is a theorem/proof draft, not a compiled
fault-tolerant implementation or an independently reviewed priority claim.
\end{abstract}

\section{Statement, conventions, and computational model}

A partition $\lambda\vdash n$ is a nonincreasing list of positive integers
with sum $n$. Its Young diagram has boxes $(r,s)$ with $1\le s\le\lambda_r$;
the content of that box is $s-r$. A standard tableau is equivalently a
path obtained by adding one box at a time. We encode the path by its sequence
of added row numbers, using a common $O(\log(n+1))$-bit width for every slot.
Let $d_\lambda$ be the number of such paths to $\lambda$, and let
$R_\lambda$ be Young's orthogonal representation, with positive
off-diagonal entries for adjacent transpositions in the usual tableau
convention. In particular its diagonal entry for $(j,j+1)$ is
$1/(\operatorname{cont}(j+1)-\operatorname{cont}(j))$.

The target convention is
\begin{equation}\label{eq:qft}
 \F_n\ket{g}=
 \sum_{\lambda\vdash n}\sum_{Q,R\in\mathcal P(\lambda)}
 \sqrt{\frac{d_\lambda}{n!}}\,
 \bra Q R_\lambda(g)\ket R\,\ket{\lambda,Q,R}.
\end{equation}
The shape is redundant with either path but may be retained as an output
register. Input permutations are encoded by canonical cycle transversals
$g=t_{i_n}^{(n)}\cdots t_{i_1}^{(1)}$, where
$t_i^{(m)}=(i,\ldots,m)$ and $1\le i_m\le m$.
The input promise subspace is the span of these valid transversal words,
of dimension $n!$. The output promise subspace is the span of valid
$(\lambda,Q,R)$ with both paths ending at $\lambda$, also of dimension
$\sum_\lambda d_\lambda^2=n!$. Norm comparisons below are restricted to
the input promise subspace. The circuit has an arbitrary unitary extension
on other physical bitstrings.

Our gate model is a fixed finite universal gate set, for example Clifford+$T$,
with reversible Boolean arithmetic, controlled operations, and synthesized
rotations charged at their elementary-gate cost. Public circuit parameters
are $n$ and a requested integer precision. Data-dependent loops and random
memory accesses are replaced by fixed scans and fixed sorting networks.
Every register has a public capacity. ``Clean error $\eta$'' means an
operator-norm bound on the input isometry with all auxiliary registers
initialized to zero, compared with the target output and zero auxiliaries.
This includes error in the auxiliary registers and permits reference-system
entanglement. It does not mean merely entrywise or basis-state error.

\begin{theorem}[Uniform clean QFT bound]\label{thm:main}
For $n\ge1$ and $0<\epsilon\le1/2$, there is a uniform finite-gate circuit
whose clean-input isometry $\widetilde V$ satisfies
\[
 \norm{\widetilde V-(\F_n\otimes\ket{0}_{\rm work})}_2
 \le\epsilon/2.
\]
Its gate count and live qubit count obey
\begin{align*}
 G&=O\!\left(n^{5/2}\poly(\log(n+1),\log(1/\epsilon))\right),\\
 W&=O\!\left(n\log(n+1)+\sqrt n\,
                  \poly(\log(n+1),\log(1/\epsilon))\right).
\end{align*}
The induced channel has diamond distance at most $\epsilon$ from the ideal
channel, including its zero work output. Coherent computation and erasure
of coefficients, balancing vectors, and rotation angles are included.
One-line permutation input instead of transversal input adds only
$O(n^2\polylog(n+1))$ gates.
\end{theorem}

The polynomial logarithmic factors are not optimized. This statement is
approximate, not an exact short Givens decomposition. It bounds gates and
space, not physical depth, connectivity overhead, or a practical crossover.
The representation-theoretic input is the Mackey embedding identity in
\cite[Definitions 5.16--5.18, Theorem 5.17]{BGO}; all new computational
ingredients and the modified streaming wrapper are proved below.

\section{The matrix primitive}

\begin{definition}[Integer-interlacing matrix]\label{def:matrix}
Let $k\ge2$ and let integers in $[0,L]$ satisfy
\[
 x_0>y_0>x_1>\cdots>y_{k-2}>x_{k-1},\qquad L\ge1.
\]
Define positive numbers
\[
 \alpha_i^2=\frac{\prod_j|x_i-y_j|}{\prod_{h\ne i}|x_i-x_h|},
 \qquad
 \beta_j^2=\frac{\prod_i|y_j-x_i|}{\prod_{h\ne j}|y_j-y_h|},
\]
and the $k\times k$ matrix
\begin{equation}\label{eq:A}
 A_{i0}=\alpha_i,\qquad
 A_{i,j+1}=\frac{\alpha_i\beta_j}{x_i-y_j}.
\end{equation}
For $k=1$ put $A=[1]$. A public power of two $K\ge k$ defines
$A'=A\oplus\Id_{K-k}$. Inactive parameter slots are masked; they are not
extra Cauchy nodes. Translation of all contents is harmless, but a bound
on coordinate bit length is necessary: span alone does not bound the
encoding size of an arbitrarily translated input.
\end{definition}

\begin{theorem}[Uniform controlled matrix implementation]\label{thm:core}
For $0<\delta\le1/2$, given $k$, the active interlacing nodes, and their validity masks in fixed
$K$-slot registers, one circuit, preserving these registers, implements
$A'$ with clean error $\delta$ using
\[
 K\poly(\log(K+1),\log(L+1),\log(1/\delta))
\]
gates and work qubits. Its bound holds in operator norm over arbitrary
superpositions of all valid parameter records at the public $K,L$.
The circuit computes its own weights and angles; no table of the $K^2$
matrix entries is formed.
An additional unchanged activation bit may control the operation, with
identity on its inactive sector and the same error and resource bounds.
\end{theorem}

\subsection{Orthogonality, conditioning, and absolute norm}

\begin{lemma}\label{lem:residue}
$A$ is orthogonal. For all its active coefficients,
\[
 L^{-1}\le\alpha_i\le1,\qquad 1\le\beta_j\le L,
 \qquad |A_{ij}|\ge L^{-2}.
\]
\end{lemma}
\begin{proof}
The rational function
\[
 f(z)=\frac{\prod_j(z-y_j)}{\prod_i(z-x_i)}
      =\sum_i\frac{\alpha_i^2}{z-x_i}
\]
has the displayed positive residues by interlacing. Its expansion at
infinity gives $\sum_i\alpha_i^2=1$. Evaluating at $y_j$ gives
$\sum_i\alpha_i^2/(x_i-y_j)=0$. Subtracting this equality for two distinct
$y$'s gives orthogonality of distinct tail columns. Differentiating at
$y_j$ gives
\begin{equation}\label{eq:betaidentity}
 \beta_j^{-2}=\sum_i\frac{\alpha_i^2}{(x_i-y_j)^2},
\end{equation}
which proves normalization of every tail column and hence orthogonality.

For $0<i<k-1$, separate the products for $\alpha_i^2$ on either side
of $x_i$. The left ratio is at least $1/(x_0-x_i)$: replace each
$y_j-x_i$, $j<i-1$, by the smaller $x_{j+1}-x_i$, and use
$y_{i-1}-x_i\ge1$. The right ratio is at least
$1/(x_i-x_{k-1})$ by the same telescoping argument. Endpoints have
only one ratio, at least $1/L$. Thus $\alpha_i\ge1/L$; the upper bound
follows from their squared sum. Integer separation and the span bound
in \eqref{eq:betaidentity} give $L^{-2}\le\beta_j^{-2}\le1$.
Equation \eqref{eq:A} gives the last assertion, including the first column.
\end{proof}

\begin{lemma}[Logarithmic absolute norm]\label{lem:absnorm}
For $B=|A'|$ entrywise and $C=\log_2K+1$, one has $\norm B_2\le C$.
\end{lemma}
\begin{proof}
Set $S=A'\diag(-1,1,\ldots,1)$. Its strict upper triangle is nonnegative
and its strict lower triangle nonpositive. This follows directly from
$x_i-y_j>0$ exactly when $i\le j$, and also holds on the identity padding.
For each binary splitting level of $[K]$, form the direct sum over sibling
intervals $I,J$ of
\[
 \begin{pmatrix}0&S_{I,J}\\-S_{J,I}&0\end{pmatrix}.
\]
Its norm is at most one, because every indicated submatrix is a compression
of an orthogonal matrix. Each off-diagonal entry of $|S|$ appears at
exactly one level. The absolute diagonal also has norm at most one.
The triangle inequality over $\log_2K+1$ pieces proves the claim.
\end{proof}

\subsection{Positive balancing with finite precision}

\begin{lemma}[Finite positive balancing]\label{lem:balance}
Put
\[
 H=\begin{pmatrix}0&B\\B^T&0\end{pmatrix},\quad
 \lambda=2C,\quad
 v_0=\mathbf1,\quad v_{t+1}=\mathbf1+Hv_t/\lambda.
\]
Choose $T$ to be the least nonnegative integer with $2^{2T}\ge8K$, so
$T=O(\log(K+1))$. Then $v=v_T=(l,r)$ satisfies
$v\ge\mathbf1$, $\norm v_2\le2\sqrt{2K}$, and
$Br\le\lambda l$, $B^Tl\le\lambda r$ componentwise.

Suppose stored weights define a nonnegative $\widehat B$ with
$|\widehat B-B|\le dB$, $d\le1/(10000K)$, and each approximate product
has component error at most $\tau\le\lambda/(10000K)$, including the
subsequent division/addition error in product units. Running the same
recurrence and rounding each iterate to dyadic words gives
\begin{equation}\label{eq:balance-error}
 e:=\norm{\widehat v-v}_2
 \le4\sqrt{2K}(d+\tau/\lambda)<0.002/\sqrt K.
\end{equation}
The stored entries lie in $[1/2,4K]$, and their weighted row/column sums
are at most
\begin{equation}\label{eq:stored-slack}
 \lambda+\frac{e(\lambda+1)}{1-e}.
\end{equation}
\end{lemma}
\begin{proof}
Write $M=H/\lambda$, so $\norm M\le1/2$ and
$v=\sum_{t=0}^{T}M^t\mathbf1$. Thus $Mv=v-\mathbf1+M^{T+1}\mathbf1\le v$,
since $\norm{M^{T+1}\mathbf1}_\infty\le2^{-T-1}\sqrt{2K}\le1/2$.
The other exact bounds follow from the geometric series.
For a nonnegative matrix, entrywise domination implies spectral-norm
domination by applying the matrix to the absolute value of a vector.
Consequently $\norm{\widehat H-H}\le dC$ and
$\norm{\widehat H/\lambda}\le(1+d)/2\le3/4$.
At each step the difference recurrence is bounded by a contraction of
factor $3/4$, plus
$\sqrt{2K}(d+\tau/\lambda)$, proving \eqref{eq:balance-error}.
Every row and column of $B$ has Euclidean norm one. Hence, for example,
\[
 B\widehat r\le\lambda l+e\mathbf1
 \le\lambda\widehat l+e(\lambda+1)\mathbf1.
\]
Divide by $\widehat l_i\ge1-e$ to obtain \eqref{eq:stored-slack}.
This uses a Euclidean error bound, not an infinity-norm bound alone.
The error and entry-range bounds also hold at every intermediate iterate,
by the same contraction induction; only the final subinvariance statement
needs the chosen value of $T$.
\end{proof}

\subsection{Fast potentials with a noncircular bit-complexity proof}

\begin{lemma}[Fixed-schedule rational potential evaluation]\label{lem:potential}
Let $N\ge1$, $Q\ge1$, and $0<\zeta\le1$. Let at most $N$ distinct integer source nodes and at most $N$ integer
targets lie in $[0,L]$. Charges are signed dyadics with a common denominator
$2^b$ and $|q_s|\le Q$. Self terms are excluded. Either potential
\[
 \sum_{s\ne t}q_s\log|t-s|,
 \qquad
 \sum_{s\ne t}\frac{q_s}{|t-s|}
\]
can be computed to absolute error $\zeta$ at all targets by a fixed
reversible schedule of size
\[
 N\poly\bigl(b,\log(N+1),\log(L+1),\log(Q+1),\log(1/\zeta)\bigr).
\]
The same bound holds for retained work. Inputs may be coherent. Padded
inactive records never participate in singular evaluations.
\end{lemma}
\begin{proof}
Let $D=\lceil\log_2(L+1)\rceil$ and use boxes of widths $w=2^\ell$,
$0\le\ell\le D$. At each level a target queries children of its parent's
box and of that parent's two neighbors, omitting its own box and its
immediate neighbors. There are at most three contributing source boxes;
six fixed query slots suffice. For a queried source box with center $c$,
put $z_s=2(s-c)/w$ and $r=w/(2(t-c))$. Then $|z_s|\le1$ and $|r|\le1/3$.
Every pair at integer distance at least two occurs at exactly one such
interaction: choose the last level at which its boxes are nonneighbors,
whose parents are therefore neighbors. Distance-one pairs are handled
exactly; their logarithmic contribution is zero. Their inverse contributions
are obtained by two additional fixed join queries with keys $t-1,t+1$.
Self pairs are omitted.

Compute moments $M_j=\sum_s q_sz_s^j$, $0\le j\le p$, in each occupied
box. The far-field formulas are the truncated series
\begin{align*}
 \sum_s q_s\log|t-s|
  &=M_0\log|t-c|-\sum_{j=1}^{p}\frac{M_jr^j}{j}+E_{\log},\\
 \sum_s\frac{q_s}{|t-s|}
  &=\frac1{|t-c|}\sum_{j=0}^{p}M_jr^j+E_{\rm inv}.
\end{align*}
Over the whole interaction list, the absolute tails are at most
$3NQ3^{-p-1}/(2(p+1))$ and $NQ3^{-p-1}$ respectively. The inverse bound
also follows from its per-source relative tail $|rz_s|^{p+1}$ and
$|t-s|\ge1$. Choose $p\ge1$ so $NQ3^{-p-1}\le\zeta/4$.

All polynomial terms are evaluated as \emph{exact rational numbers}.
To approximate a scalar logarithm, range-reduce its positive rational
argument to $2^a u$ with $1\le u\le2$ and use
\[
 \log u=2\sum_{j=0}^{R-1}\frac{z^{2j+1}}{2j+1}+E_R,
 \quad z=\frac{u-1}{u+1},\quad |E_R|\le3^{-2R}.
\]
Approximate $\log2$ by the same formula at $z=1/3$.
For the half-integer distances here, the scalar error is at most
$(D+5)3^{-2R}$. Choose $R$ so
$NQ(D+5)3^{-2R}\le\zeta/4$. Interactions partition the sources, so this
is also the aggregate error bound. Round the final potential once to a
dyadic with error at most $\zeta/2$.

The following denominator accounting is essential. At level $\ell$ the
$j$th moment has a common denominator dividing $2^{b+\ell j}$.
Its numerator therefore has
$O(b+pD+\log(N(Q+1)))$ bits, \emph{not} $N$ times that many bits.
For a far logarithmic polynomial a common denominator may include
$2^{b+pD}|2(t-c)|^p p!$, with polynomial bit length in $b,p,D$;
the inverse polynomial can require the additional power $|2(t-c)|$.
The $R$-term logarithm has a denominator of bit length
$O(R(D+\log(R+1)))$, using a common power of its argument denominator
and $(2R-1)!$. A target sums only $O(D)$ interaction values. Thus their
denominators can be multiplied even without cancellation while retaining
polynomial bit length. An explicit loose envelope is obtained by setting
\begin{align*}
 A_0&=\lceil\log_2(NQ+1)\rceil,\\
 W_{\log}&=4R\bigl(D+4+\lceil\log_2(2R+1)\rceil\bigr)
             +\lceil\log_2(D+5)\rceil+8,\\
 W_0&=b+(2p+2)(D+4)+p\lceil\log_2(p+1)\rceil\\
    &\hspace{1cm}+W_{\log}+A_0+\lceil\log_2(D+5)\rceil+16.
\end{align*}
The indicated common denominators and the numerator magnitude bounds fit
each far rational in $W_0$ bits. Summing the at most $6D+2$ nontrivial
padded slots and rounding to $b_{\rm out}$ fractional bits fits within
$4(6D+3)(W_0+b_{\rm out}+1)$ bits. The unused top level may be skipped:
all nodes occupy a single width-$2^D$ box. Integer arithmetic
on these words has polynomial Boolean cost. There is no stability
assumption based merely on bounded intermediate magnitudes and no
precision chosen in terms of itself.

For access, sort fixed padded source records by node and use segmented
scans to compute moments. One representative per occupied box is marked
as a source record. Form the six queries per target at fixed addresses.
A fixed sorting network on their box keys joins source payloads to queries;
a scan broadcasts the appropriate moments within equal-key runs. Keep
comparison decisions and scan history. Reverse the sorting, copy answers
from the original fixed query slots, and reverse all the work. This costs
$O(N\log^2(N+1))$ comparisons per level, with polynomial-width payloads.
Moment reductions use the known common dyadic denominator, rather than
naive repeated multiplication of rational denominators. A bounded Boolean
computation is made reversible by recording its steps, copying its result,
and reversing it. These constructions give the asserted gate and space
bound, including the data movement.
\end{proof}

\begin{lemma}[Coefficients, balancing, and scalar functions]\label{lem:arithmetic}
All classical data required by Lemma~\ref{lem:balance} can be computed,
stored as dyadics, and erased with
$K\poly(\log(K+1),\log(L+1),\log(1/\eta))$ gates and work, at any
requested relative weight accuracy $0<\eta\le1$ and the stated balancing
accuracy. The scalar functions required for probability trees have
polynomial bit cost at this precision.
\end{lemma}
\begin{proof}
Put charge $-1$ at every $x_i$ and $+1$ at every $y_j$. The self-excluded
log potential equals $\log\alpha_i^2$ at $x_i$ and
$-\log\beta_j^2$ at $y_j$. Lemma~\ref{lem:potential} computes these with
absolute error $O(\eta)$, which implies relative error $O(\eta)$ after
exponentiation and a square root, or directly after halving the logarithm.
Lemma~\ref{lem:residue} bounds the true logarithms by $2\log L$ in
absolute value. Truncation to $b=\poly(\log(KL/\eta))$ dyadic bits preserves
the requested relative error because the weights are polynomially bounded
above and below.

For completeness, put $D=\lceil\log_2(L+1)\rceil$ and let the dyadic
half-logarithm argument $a$ satisfy $|a|\le D+2$. Choose a public power
of two $h$ with $2(D+2)\le h<4(D+2)$. Approximate $\exp(a/h)$ by its
rational Taylor polynomial of degree $R_e$, then square $\log_2h$ times.
The Taylor tail is at most $3\cdot2^{-R_e}$, and both the exact and
approximated reduced exponentials are positive and at most two.
The inequality $|u^h-v^h|\le h2^{h-1}|u-v|$ shows that
\[
 R_e\ge h+\left\lceil\log_2(4h/\zeta_e)\right\rceil
\]
suffices for absolute error $\zeta_e$, with slack for final rounding.
For a $b_a$-bit argument the rational word length is
$O(hR_e(b_a+\log(h+1)+\log(R_e+1)))$, hence polynomial in the logarithmic
parameters. Computing the log potentials to $\eta/8$ and their signed
half-exponentials to $\zeta_e=\eta/(4L)$ yields relative weight error
less than $\eta$. Round the result to a dyadic before further large sums.
Integer square
roots give a square root to any dyadic precision by scaling and integer
division. Thus no exact-real primitive has been assumed.

For matrix-vector multiplication the active block satisfies
\[
 (Bz)_i=\alpha_i z_0+
       \alpha_i\sum_j\frac{\beta_jz_{j+1}}{|x_i-y_j|}.
\]
For $B^Tz$, its first entry is $\sum_i\alpha_i z_i$ and its remaining
entries are $\beta_j\sum_i\alpha_i z_i/|y_j-x_i|$.
Use the stored dyadic weights and vectors as exact inputs to
Lemma~\ref{lem:potential}; their charges have a common dyadic denominator.
Multiplication by the outside weights and the constant-column sums have
polynomial bit cost. Identity-padded indices are handled analytically.
At each iterate the induction bound in Lemma~\ref{lem:balance} gives
charges at most $Q=8KL$, since stored weights are at most twice their
exact upper bounds. A potential tolerance $\tau/(16L)$, with sufficiently
fine final scalar rounding, attains the prescribed $\tau$. If each weight
has relative error at most $u\le1$, then the resulting factored matrix has
entrywise relative error at most $3u$, so take $u\le d/3$ for balancing.
Round \emph{each iterate}
back to dyadics. This prevents denominator growth across iterations.
There are $O(\log(K+1))$ iterations, so the total remains near-linear.

For positive sibling masses, put $r=\min(L_0,R_0)/\max(L_0,R_0)$ and
$u=\sqrt r\in[0,1]$. The tree angle is $4\arctan w$, or
$\pi-4\arctan w$, with $w=u/(1+\sqrt{1+u^2})<1/2$.
This avoids taking the reciprocal of a small approximate number.
Round each square root and the division to absolute error at most $h_0$.
Since $|w'(u)|\le1/2$ and the denominator is bounded below by two before
rounding, the resulting $w$ error is at most $2h_0$.
For $h_0\le1/32$ it remains below $1/2$.
An $R_0'$-term rational alternating arctangent series then gives angle
error at most $9h_0+2^{-2R_0'}$, apart from the public $\pi$ error.
Taking $h_0\le\epsilon_\theta/32$ and each remaining error at most
$\epsilon_\theta/4$ suffices. The polynomial's common denominator has
$O(R_0'(b+\log(R_0'+1)))$ bits for a $b$-bit dyadic argument.
An integer square root of
$\lfloor2^{2b}\,\text{numerator}/\text{denominator}\rfloor$
produces a dyadic square-root approximation with error less than $2^{-b}$.
Public $\pi$ can be computed, for
example, by $\pi=16\arctan(1/5)-4\arctan(1/239)$.
Sine uses its factorially convergent Taylor series. The public
normalization below is the reciprocal of a sine of inverse-logarithmic
size. All these calculations use fixed polynomial-length integer routines
and fixed truncation lengths chosen from the desired accuracy.
\end{proof}

\subsection{Clean overlap amplification}

Fix the stored positive vectors $\widehat l,\widehat r$ from balancing.
Set the public integer $m=2\lambda$, and
\[
 \theta=\frac{\pi}{4m+2},\qquad \mu=\csc\theta.
\]
Elementary sine bounds give $2\lambda\le\mu\le4\lambda+1$.
In view of \eqref{eq:stored-slack}, the following padding probabilities
are greater than $0.49$.

Let the ambient register have mutually orthogonal basis sectors
$\ket{i,j,g}$, $\ket{i,\mathrm{rowpad}}$, and
$\ket{j,\mathrm{colpad}}$. Define isometries
\begin{align}
 P\ket i&=\sum_j\sqrt{\frac{B_{ij}\widehat r_j}{\mu\widehat l_i}}
     \ket{i,j,g}+\sqrt{1-s_i}\ket{i,\mathrm{rowpad}},\label{eq:P}\\
 Q\ket j&=\sum_i\sgn(A'_{ij})
   \sqrt{\frac{B_{ij}\widehat l_i}{\mu\widehat r_j}}
     \ket{i,j,g}+\sqrt{1-t_j}\ket{j,\mathrm{colpad}},\label{eq:Q}
\end{align}
where $s_i,t_j$ are the sums of the displayed good probabilities.
The comparison isometries use \emph{exact} $A'$ and the \emph{same stored}
balancing vectors. Consequently
\begin{equation}\label{eq:overlap}
 P^\dagger Q=A'/\mu
\end{equation}
exactly; no approximate nonorthogonal matrix is being amplified as if it
were orthogonal.

\begin{lemma}[Amplification and complete-output error]\label{lem:quantum}
Let $R_P=2PP^\dagger-\Id$ and $R_Q=2QQ^\dagger-\Id$. Then
\[
 (-R_QR_P)^m Q=PA'.
\]
Suppose actual preparation unitaries $\widetilde U_P,\widetilde U_Q$
have clean-input isometries $\widetilde P,\widetilde Q$ within $\xi$ of
$P,Q$. Reflections about their \emph{complete} input-clean subspace,
followed by the actual inverse $\widetilde U_P^\dagger$, implement $A'$
with complete-output error at most $(8m+2)\xi$.
\end{lemma}
\begin{proof}
Put $X=PA'$ and $Y=(Q-\sin\theta X)/\cos\theta$.
Equation~\eqref{eq:overlap} shows that $X,Y$ are isometries with orthogonal
ranges, and $Q=\sin\theta X+\cos\theta Y$. On these paired subspaces,
\[
 R_P=\begin{pmatrix}1&0\\0&-1\end{pmatrix},\qquad
 -R_QR_P=\begin{pmatrix}\cos2\theta&\sin2\theta\\
                              -\sin2\theta&\cos2\theta\end{pmatrix}.
\]
After $m$ iterations the coefficients of $X,Y$ are respectively
$\sin((2m+1)\theta)$ and $\cos((2m+1)\theta)$. The chosen angle makes
these one and zero, proving the identity with its sign.

Let $J$ denote the logical input embedding with every target, tag,
arithmetic, and synthesis auxiliary zero. A preparation implements
$\widetilde P=\widetilde U_PJ$ (and similarly for $Q$, whose preparation
also swaps the logical input into its column-source position).
The realized reflection is
$\widetilde U_P(2JJ^\dagger-\Id)\widetilde U_P^\dagger$.
It must test \emph{all} input-clean auxiliary registers, not just a
success flag. Read-only coefficient tables are external parameters, not
preparation-clean work. For isometries,
\[
 \norm{\widetilde P\widetilde P^\dagger-PP^\dagger}\le2\xi,
 \qquad\norm{\widetilde R_P-R_P}\le4\xi.
\]
Telescoping the $2m$ reflections and initial preparation bounds the
amplified-isometry error by $(8m+1)\xi$.
The final inverse requires no closeness assertion about arbitrary full
unitary extensions: since $\widetilde U_P^\dagger\widetilde P=J$,
\[
 \norm{\widetilde U_P^\dagger\widetilde W-JA'}
 \le\norm{\widetilde W-PA'}+
       \norm{(P-\widetilde P)A'}\le(8m+2)\xi.
\]
The minus sign in the amplification identity is retained as a controlled
phase whenever the whole operation is conditionally applied.
\end{proof}

\begin{lemma}[Preparation cost and error]\label{lem:prep}
For $0<\xi\le1/2$, the preparations in \eqref{eq:P}--\eqref{eq:Q} admit the preceding
approximations with cost
$K\poly(\log(K+1),\log(L+1),\log(1/\xi))$ per preparation or inverse.
\end{lemma}
\begin{proof}
For a coherent source index, one fixed scan retrieves its constant number
of coefficient words. Visit each target record once at its fixed address
to compute the $K$ good probabilities, followed by the padding leaf.
Use a fixed $2K$-leaf tree: $K$ good leaves, one private-padding leaf, and
$K-1$ exact-zero leaves. A tag bit and target bits represent these leaves;
a fixed reversible relabeling gives the row-private or column-private tag.
There is \emph{one} length-$K$ list controlled by the source, not $K$
lists for its $K$ possible basis values. Build its binary sum tree in
$O(K)$ word operations; perform one prefix-controlled rotation for each
of its $O(K)$ internal nodes. The source is unchanged, so all coefficient,
sum, and angle arithmetic can be reversed after the target is superposed.
For $Q$, compute the negative-entry predicate on the good sector, apply
$(-1)^{[A'_{ij}<0]}$, and erase the sign scratch. At zero entries this
phase is $+1$, so it remains unitary even on leakage states.
All lookups are fixed scans or fixed addresses.

Every nonzero good leaf is at least
\[
 p_{\min}=\frac1{8K(4\lambda+1)L^2},
\]
by Lemma~\ref{lem:residue} and the stored-vector bounds. Exact zero leaves
from identity padding or inactive slots are identified logically and stay
zero. Compute each nonzero leaf with relative error at most $\eta$,
rounding it to a \emph{common dyadic denominator before summing}.
An absolute rounding target $O(\eta p_{\min})$ suffices; it uses only
logarithmic additional bits. Define the padding as one minus this dyadic
sum. Its constant slack implies relative error $O(\eta)$ there too.
All positive subtree sums consequently have relative error $O(\eta)$.

A tree rotation has angle $2\arctan\sqrt{R/L}$ for left/right masses
$L,R$. Its derivative with respect to $\log(R/L)$ has absolute value at
most $1/2$. Thus a uniform relative mass error produces $O(\eta)$ angle
error even for unbalanced branches. Zero-child cases use exact endpoint
rotations, bypassing division by zero; an all-zero subtree has angle zero.
More explicitly, let $h=2K-1$. A leaf relative error at most $\eta$ on
good leaves induces relative error at most $2\eta$ on the padding and
subtrees, and angle error at most $4\eta$, for $\eta\le1/4$.
If angle evaluation error is $a$ and the full controlled-rotation synthesis
error is $s$, the preparation error is at most $h(2\eta+a/2+s)$.
Choose $\eta=\xi/(16h)$ and $a=s=\xi/(4h)$, leaving slack below $\xi$.
For a superposed source, each error is a direct sum of
source-controlled blocks and has their maximum norm, not their sum.

A binary angle word controls a fixed sequence of public binary-angle
rotations. Each public rotation can be compiled over the chosen finite
gate set with polylogarithmic precision overhead, for example by
\cite{DN}. Controlled rotations are synthesized with their full phase
convention, not just up to independent source-dependent phases. For example,
with $V=R_y(\gamma/2)$, use
\[
 \operatorname{controlled}\!\text{-}R_y(\gamma)
 =(\Id\otimes V)\,\mathrm{CNOT}\,
  (\Id\otimes V^\dagger)\,\mathrm{CNOT}.
\]
Replacing $V$ by a synthesized $\widetilde V$ and using its exact circuit
inverse cancels its global phase; the inactive controlled block is exactly
identity and the active error is at most $2\norm{\widetilde V-V}$.
An external activation control can be included in each reversible predicate
and rotation control with the same asymptotic overhead. Its synthesis
errors are included in the same budget, and inactive behavior is identity.
The complete-clean reflection tests a number of bits no greater than
the work bound, so it has the same asymptotic cost. This includes the
reversible tape for the preparation arithmetic.
\end{proof}

\begin{proof}[Proof of Theorem~\ref{thm:core}]
Compute dyadic weights sufficiently accurately both for balancing and for
Lemma~\ref{lem:prep}. Lemmas~\ref{lem:potential}--\ref{lem:arithmetic}
give near-linear reversible preprocessing, including the $O(\log K)$
balancing iterations. Take $\xi=\delta/(8m+2)$ and use
Lemmas~\ref{lem:quantum}--\ref{lem:prep}. There are $O(\log K)$
preparations/reflections. After the final inverse, reverse all
parameter-only preprocessing, since the parameter registers never changed.

All functions are totalized on invalid encodings using validity controls
and safe dummy denominators. Every circuit gate is therefore unitary on
the full register space, although the accuracy promise is only on valid
inputs. One fixed schedule and one public $m,\mu$ serve every valid $k$
and parameter record. Their orthogonal control sectors imply the same
operator-norm error bound for arbitrary superpositions, including phases.
No enumeration or separate compilation of shape sectors is required.
\end{proof}

\section{A compact streaming wrapper for the symmetric group}

Throughout this section fix global public capacities
\[
 D=\lfloor\sqrt n\rfloor,\qquad
 K=2^{\lceil\log_2(2D+2)\rceil},\qquad L=2n+1.
\]
The same capacities are used at every outer and inner stage. Contents of
any prefix can be shifted by $n$ into $[0,L]$ without changing its matrix.

\begin{lemma}[Compact shape operations]\label{lem:shape}
A partition with at most $n$ boxes has a canonical $O(D)$-word encoding
supporting valid box addition/removal, row-length queries, corner
extraction, and reversible row/corner relabeling with
$D\polylog(n+1)$ gates and work per operation. A word has
$O(\log(n+1))$ bits. No evolving list of previous shapes is required.
\end{lemma}
\begin{proof}
If there are $d$ diagonal boxes, store $d$ and the Frobenius coordinates
$a_i=\lambda_i-i$, $b_i=\lambda'_i-i$, $1\le i\le d$, with zero inactive
slots. Since $d^2\le n$, there are at most $D$ active slots. Row length is
\[
 \lambda_r=\begin{cases}
 a_r+r,&r\le d,\\
 \sum_{i=1}^{d}[b_i+i\ge r],&r>d.
 \end{cases}
\]
A fixed scan implements both cases. Adding a valid box $(r,s)$ increments
$a_r$ when $s>r$, increments $b_s$ when $s<r$, and appends $(0,0)$ and
increments $d$ when $s=r$. Removal reverses this. For a clean in-place
update with retained row $r$, compute the new canonical encoding into
zero work, erase calculation scratch, swap old/new, and erase the old
encoding by the inverse removal function. All indexing uses fixed scans.

Writing $\delta_c$ for a point mass at content $c$, the corner identity is
\[
 \delta_0+\sum_{i=1}^{d}
 (\delta_{a_i+1}-\delta_{a_i}+\delta_{-b_i-1}-\delta_{-b_i})
 =\sum_{x\in\AC(\lambda)}\delta_x-
   \sum_{y\in\RC(\lambda)}\delta_y.
\]
It follows either by following the diagram boundary or by canceling
successive horizontal and vertical boundary edges. Sorting and aggregating
these $O(D)$ records yields the alternating corner contents and their
validity masks. The number of addable corners is at most $2d+1$, hence
fits $K$. For an explicit row-label construction, use the following
fixed masked candidates, with each pair denoting (row, content).
For $1\le i\le d$, upper addable candidates are $(i,a_i+1)$ when
$i=1$ or $a_{i-1}\ge a_i+2$; lower addable candidates are
$(i+b_i+1,-b_i-1)$ when $i=1$ or $b_{i-1}\ge b_i+2$.
Also include $(d+1,0)$ if $d=0$ or both $a_d,b_d>0$.
Upper removable candidates are $(i,a_i)$ when $a_i>0$ and
either $i=d$ or $a_i\ge a_{i+1}+2$; lower removable candidates are
$(i+b_i,-b_i)$ when $b_i>0$ and either $i=d$ or
$b_i\ge b_{i+1}+2$. Include $(d,0)$ if $d>0$ and $a_d=b_d=0$.
The upper cases translate the neighboring-row inequalities for an outer
or inner corner; the lower cases follow by transposition. The diagonal
cases check the two boundary edges adjacent to the diagonal.
Each candidate list has at most $2D+1$ public slots, including inactive
ones. Sort the valid candidates by decreasing content.
Fixed sorting, copying, and reverse sorting erase the comparison history.
A scan converts a row to its corner ordinal, and an inverse scan erases
the old label. Corners are ordered by descending content; the first
column label is the separate star symbol, followed by removable corners.
\end{proof}

\subsection{The imported algebraic identity}

For $\lambda\vdash M-1$, define the induction isometry
\begin{equation}\label{eq:ind}
 U_{M,\lambda}\ket{i,P}=
 \sum_{\mu\succ\lambda}\sum_{Q\in\mathcal P(\mu)}
 \sqrt{\frac{d_\mu}{M d_\lambda}}\,
 \bra Q R_\mu(t_i^{(M)})\ket{P\to\mu}\,\ket{\mu,Q}.
\end{equation}
We use the established Mackey recurrence \cite[Eq.~(66)]{BGO}:
after recursively inducing the truncated path, the final local block
is $A_\lambda$ when its two intermediate shapes coincide, and is the
scalar $+1$ on the unique detour when they differ. The identity-transversal
branch is the star column of $A_\lambda$. To identify it with
Definition~\ref{def:matrix}, the branching dimension ratio in
\cite[Lemma 2.4]{BGO} gives
\[
 \alpha_i^2=\frac{d_{\lambda+\square_i}}{M d_\lambda}.
\]
Normalizing the removable column in \cite[Eq.~(64)]{BGO} and using
\eqref{eq:betaidentity} then identifies its positive scale as
$\beta_j^2=(M-1)d_{\lambda-\square_j}/d_\lambda$.
Its denominator is exactly the difference of the addable and removable
contents, so the signed entries agree as well. The extended
map for $i>M$ returns the unchanged endpoint, path $(P,\star)$, and flag
$i-M+1$; valid inputs return flag $1$. These are algebraic inputs from
\cite{BGO}, not consequences of numerical tests or new claims here.

\begin{proposition}[Streaming induction]\label{prop:stream}
Assume a clean, uniformly shape-controlled implementation of $A_\lambda$
at the above public capacity. Then $U_{M,\lambda}$ can be implemented by
$M$ local matrix calls and $MD\polylog(n+1)$ additional gates, including
compact controls and all boundary cleanup.
\end{proposition}
\begin{proof}
Write the original tableau word as $p_1,\ldots,p_{M-1}$ and its original
prefix shapes as $P_j$. Its endpoint $\lambda=P_{M-1}$ is an external
control. Initialize an internal word
$(\star,p_1,\ldots,p_{M-1})$, compact prefix $\varnothing$, token $a=0$,
and flag $f=i$. The value zero encodes the star, not a physical row.
The word is the existing left path relabeled together with one fresh zero
slot, not a retained copy of the input path. The transversal input register
itself serves as $f$, so there is no leftover copy of $i$ at the boundary.

At stage $j>1$ first add the untouched original row $p=p_{j-1}$ to the
compact prefix. If $a=0$, swap the preceding star slot with the current
word slot and decrement $f$. At $j=1$ omit the prefix update, swap,
and flag decrement.
If $f\ne1$, do nothing further. Otherwise the local embedding is as
follows. Compute the equality bit $e=[a=p]$, with $p$ denoting the current
input symbol, which can be star at activation. If $e=0$, leave $p,a$
unchanged. If $e=1$, erase the duplicate token by XOR with $p$, convert
$p$ from star/removable-row to the corresponding column index, apply
$A_{P_{j-1}}$, convert the output index to addable row $q$, and copy $q$
into the token. Erase $e$ using the output predicate $[a=q]$.

Here is the complete induction invariant after stage $j$. The compact
prefix is the \emph{original} $P_{j-1}$ and the suffix $p_j,\ldots,p_{M-1}$
is untouched. If $i>j$, the processed word is
$(p_1,\ldots,p_{j-1},\star)$, token zero, and flag $i-j+1$.
If $i\le j$, its amplitudes and processed path equal
$U_{j,P_{j-1}}$ on the original prefix input; the endpoint is
$P_{j-1}+\square_a$, token $a>0$, and flag one.

The empty-shape matrix proves the base case. For an inactive branch the
swap moves the star and the flag as claimed; when $i=j$ the star column
activates the identity-transversal branch. For a previously active branch,
let $\kappa=P_{j-2}$. The new original endpoint is
$\kappa+\square_p$ and the previously induced endpoint is
$\kappa+\square_a$. They coincide exactly when $p=a$.
When distinct, their common successor is
$\kappa+\square_p+\square_a$, so the imported detour block $+1$ outputs
$q=p$ and retains $a$. When equal, the imported block is precisely the
matrix just applied. This proves both the amplitudes and endpoint invariant.
It also proves that $[a=q]$ is the output equality predicate in both cases.
The token is not generally the final output path symbol.

Each step is reversible without an edit log. If $f=1$ and the output is
on the diagonal $a=q$, erase the duplicate, apply $A^T$, and copy the recovered input
back to the token. A recovered zero token identifies the former inactive
branch, whose flag increment and swap can be undone for $j>1$. Then remove
the restored original row from the compact prefix, also only for $j>1$.
The compact control is
unchanged during every matrix call, including its preprocessing erasure.

For valid inputs, at the boundary copy the token $a$ to the new right-path row slot before
erasing it, and form $\mu=\lambda+\square_a$. Define the exact checksum
\[
 W(\lambda)=\sum_r r\lambda_r
 =\sum_{i=1}^{d}\left(i(1+a_i+b_i)+\frac{b_i(b_i+1)}2\right).
\]
Its Frobenius formula partitions the boxes into diagonal hooks.
Under the one-box promise $a=W(\mu)-W(\lambda)$, so the internal token
can be erased. The checksum needs $O(\log(n+1))$ bits with enough range
for $n(n+1)/2$; it is not a general-purpose shape hash.
Erase the prefix against external $\lambda$, and erase the valid flag one.
In the full QFT erase the old external $\lambda$ by reconstructing it
from $\mu$ and the retained right-path row; promote $\mu$ to the endpoint.
The invalid extension bypasses box addition and checksum erasure and
instead retains its specified star/flag output.
Every stage uses a fixed number of compact scans and relabelings, so
Lemma~\ref{lem:shape} gives the claimed cost. If expanded endpoint lists
are desired at an outer boundary, even $O(MD)$ conversion work is within
the same bound; none is done inside the inner loop.
\end{proof}

\subsection{Completion of the Fourier transform and resource sum}

\begin{proof}[Proof of Theorem~\ref{thm:main}]
First consider exact local matrices. Write $g=t_i^{(M)}h$ with
$h\in S_{M-1}$. Suppose the previous transform has left/right paths
$P,R$ of shape $\lambda$ and coefficient
$\sqrt{d_\lambda/(M-1)!}\,[R_\lambda(h)]^P_R$.
Apply \eqref{eq:ind} to $(i,P)$ and append the box $\mu\setminus\lambda$
to the \emph{right} path. For any fixed output right path
$S=R\to\mu$, the penultimate shape $\lambda$ is unique, and the output
coefficient is
\[
 \sqrt{\frac{d_\mu}{M!}}
 \sum_P [R_\mu(t_i)]^Q_{P\to\mu}[R_\lambda(h)]^P_R
 =\sqrt{\frac{d_\mu}{M!}}[R_\mu(g)]^Q_S.
\]
The equality is subgroup adaptation of Young's orthogonal representation.
Starting at $M=1$ proves exactly the convention \eqref{eq:qft}.
Proposition~\ref{prop:stream} supplies the clean implementation, rather
than discarding intermediate endpoint registers.

Use Theorem~\ref{thm:core} at common local error
$\delta=\epsilon/[2n(n+1)]$. There are
$N=n(n+1)/2$ calls, each at $K=O(\sqrt n)$ and $L=O(n)$.
All other word operations are exact reversible Boolean circuits.
Telescoping uses approximate suffixes and \emph{ideal prefixes}, so every
local error is applied on its promised valid subspace even if an actual
earlier approximation leaks outside it. The resulting clean isometry
error is at most $N\delta=\epsilon/4$, in particular at most $\epsilon/2$.
The diamond distance of the associated isometric channels is at most twice
their operator-norm distance; partial trace cannot increase it.

The total cost of calls and wrappers is
\[
 O\!\left(n^2\sqrt n\,
   \poly(\log(n+1),\log(1/\epsilon))\right).
\]
Two full paths and the transversal registers use $O(n\log(n+1))$ bits.
Only one local matrix workspace, compact prefix, and token/flag are live
at a time. This yields the stated qubit bound, even with the elementary
record--copy--reverse simulation of all local arithmetic.

For one-line input, at size $M$ save $i=g(M)$, apply $(t_i^{(M)})^{-1}$
to each stored image, and clear the resulting constant $M$ in the last
position. This maps image $i$ to $M$, images greater than $i$ down by one,
and smaller images to themselves. Recurse on the first $M-1$ positions.
Fixed scans cost $O(n^2\polylog(n+1))$ gates; reversing them reconstructs
the original permutation, so the conversion is clean.
\end{proof}

\section{Scope, verification, and what is still not established}

The proof is a constructive asymptotic circuit-existence argument. Its
essential new dependencies are Lemmas~\ref{lem:absnorm},
\ref{lem:balance}, and \ref{lem:potential}, the fully accounted
preparations, and Proposition~\ref{prop:stream}. Weighted square-root
embeddings and amplification have precedents in \cite{Childs,Berry};
we have given the needed identities and robustness argument explicitly.
The rational potential proof intentionally replaces an earlier informal
global floating-point rounding estimate. It does not claim that the
current interval-arithmetic prototype executes that exact rational
algorithm or proves its uniform bound by testing.

The repository contains independent small-size checks of residue matrices,
compact shape updates, intermediate streaming induction, whole small QFTs,
reversible sorted joins, numerical enclosures, and amplification. Added
theorem-audit checks specifically examine full-workspace reflections,
unrelated unitary extensions, coherent phases, outer endpoint cleanup,
and transversal conversion. These can refute implementation mistakes;
they cannot establish the theorem for arbitrary size.

The following are not supplied: a complete compiled Clifford+$T$ circuit
for the whole construction, tight logarithmic factors or constants, a
practical resource estimate, an exact near-linear Givens factorization,
an optimality lower bound, or independent expert peer review. The
mathematical claim above is a research theorem draft to audit, not an
assertion that those subsequent milestones have already been met.
Its development and internal checking were AI-assisted. Attribution of
priority beyond the explicit cited ingredients requires a separate,
systematic literature review.

\begin{thebibliography}{9}
\bibitem{BGO}
C. Bruinsma, D. Grinko, and M. Ozols,
\emph{Quantum Fourier transform for the symmetric group},
arXiv:2608.28569v2 (2026).
\href{https://arxiv.org/html/2608.28569v2}{Full text}.

\bibitem{Childs}
A. M. Childs,
\emph{On the relationship between continuous- and discrete-time quantum walk},
arXiv:0810.0312.
\href{https://arxiv.org/abs/0810.0312}{arXiv record}.

\bibitem{Berry}
D. W. Berry, A. M. Childs, R. Cleve, R. Kothari, and R. D. Somma,
\emph{Exponential improvement in precision for simulating sparse Hamiltonians},
arXiv:1312.1414.
\href{https://arxiv.org/abs/1312.1414}{arXiv record}.

\bibitem{DN}
C. M. Dawson and M. A. Nielsen,
\emph{The Solovay--Kitaev algorithm}, arXiv:quant-ph/0505030.
\href{https://arxiv.org/abs/quant-ph/0505030}{arXiv record}.
\end{thebibliography}
\end{document}
